\subsection*{先确定接收身份和效果身份}

一项付款成功可以同时供履约、审计和查询投影使用。履约建单以后，审计也有理由登记这笔付款。若全系统只存一列 \code{event\_id}，履约抢先插入 EV1，审计就可能把自己的合法工作当成重复。

给接收记录增加稳定的处理语义范围。本章履约处理器的 \code{handler\_scope} 为 \code{fulfillment.\allowbreak initial.v1}，含义是按照既定规则处理初次履约，同一处理器的多个实例共享它。它也标识这项业务订阅的消费责任。连接名、实例名、容器编号和程序构建版本会随运行变化，处理语义范围在这些变化中保持。审计使用另一范围 \code{audit.paid.v1}，保存自己的业务效果。为讨论这个独立场景，设审计从授权的历史导出接口取得 EV1 完整原件。第九章的队列目标只有 \code{fulfillment.paid}，其中多个消费者竞争取得消息。历史导出接口则向审计提供一份独立输入，审计依据自己的接收记录继续。

事件号还具有发布者的命名范围。本章只接收既定购买服务经授权通道发布的消息，固定其来源为 \code{purchase-service}。若以后另一发布者也能使用 EV1，就以来源列区分二者。接收入口根据通道身份、订阅权限和来源核验结果填写这列，随后才解释来件中的业务字段。

新路径的接收身份记为
\[
I=(\text{处理语义},\text{可信来源},\text{消息种类},\text{消息号}),
\]
示例为 \code{(fulfillment.initial.v1,purchase-service,event,EV1)}。增加消息种类是为了稍后同时容纳没有 \code{event\_id} 的旧通知，避免旧编号与新事件号相撞。Idempotent Consumer 模式原文用 \code{(subscriberId,messageID)} 作为联合主键，并在同一事务中插入消息记录与修改业务~\cite{idempotent-consumer,handling-duplicate-messages}。这里把来源与旧格式范围进一步写明。

事件唯一还留下另一个问题：两个不同事件都报告 O1 已付款时，消费者应建立几张初次履约单？为这项业务动作另设效果身份
\[
B=(\text{order\_id},\text{action},\text{authorization\_id}),
\]
O1 的值为 \code{(O1,initial,initial:O1)}。\code{initial:O1} 是业务规则为这份已付款订单确定的初次履约授权标识。消费者根据既定事件种类和可信事实推导它，将同一 O 的初次履约集中到这个固定授权上。

于是两个层次可以同时成立：I1、I2 是不同接收身份，B1 却相同，两份接收结果都可以引用 F1。Alice 用 K2 合法购买同样商品、同样金额时，得到另一 O2 和 \code{initial:O2}，对应自己的 F2。

\CTenFigure{ch10-original02}{完整事件绑定与初次履约效果绑定分属两个层次}{fig:ch10-input-bindings}
\CTenFigure{ch10-identity-map}{接收身份随处理范围与消息变化，初次履约按固定业务授权汇合}{fig:ch10-identity-map}

\subsection*{重复时比较保存的输入}

处理器先核验通道来源、订阅权限、报文结构与字段类型，再占用 I。EV1 继续使用第九章的完整报文：事件种类 \code{order.payment\_succeeded}，对象 O1，\code{schema\_version=1}，原 \code{occurred\_at}，以及保存的 P1/C1/O1、账户、商品、数量、报价、币种金额和平台范围/X1。消息对象与 \code{payload.order\_id} 必须一致。数量和最小货币单位金额采用正整数，版本字段也按整数类型核验。例如 \code{schema\_version=true} 在类型检查时进入结构错误分支，保存完整来件与核查责任。

接收记录绑定完整语义报文。JSON 对象成员的排列和传输空白不影响比较，字段的类型和值、数组的顺序有含义。这里按保存格式解释并比较结构化值，整数 1 与布尔值 \code{true} 分别处理。PA1、delivery tag 和本次链路号保存在传输与尝试记录中，可以随发送变化。接收绑定比较的是保存的事件正文，这些传输变化各有自己的记录位置。

效果输入取出原 \code{payload} 中与本次建单有关的全部字段：P/C/O、账户、商品、数量、Q7、CNY 3000、原平台范围和 X1。EV1、事件记录时间和发布尝试号用于定位事件与发送。这样，新补建事件与旧通知可以拥有不同完整输入，同时报告同一项已付款初次履约授权。

购买库的商品现价后来改为 Q8 的 3500 分，处理器按 EV1 保存的 Q7 和 3000 分识别原成交。第一次接收以后再读当前价格拼输入，会使一次正常重投变成异参，甚至按新价格生成另一项效果。需要当前库存、禁运或调度政策时，把它们作为首次业务判断的依据，保存当时结论。已有 I 的恢复先于重新作这类判断。

\subsection*{第一层：接收记录、授权与实际履约}

先看决定初次履约的三张核心表。\code{authorizations} 保存 B 及其效果输入，\code{fulfillments} 保存实际 F，\code{processed\_events} 保存 I 的完整输入与本次结果。正常成功沿 I 找到 B，再关联真实 F。两层唯一约束分别约束接收与效果，外键检查同库的授权归属。

下面的 DDL 直接给出本章最终使用的列，避免逐段切换不同版本的表。首次运行先看 \code{allowed/rejected} 与来源 \code{consumer}。\code{hold} 是暂停自动建单的授权决定，\code{legacy-import} 标记有依据的旧履约导入；第~\ref{section:10-3}~节在升级问题出现时完整推导这两项取值及 \code{origin\_evidence} 的来源。

\begin{BookCode}
CREATE TABLE authorizations (
  order_id text NOT NULL, action text NOT NULL CHECK(action='initial'),
  authorization_id text NOT NULL, effect_input jsonb NOT NULL,
  disposition text NOT NULL CHECK(disposition IN ('allowed','rejected','hold')),
  reason text, PRIMARY KEY(order_id,action,authorization_id),
  CHECK(authorization_id='initial:'||order_id)
);
CREATE TABLE fulfillments (
  fulfillment_id bigint GENERATED BY DEFAULT AS IDENTITY PRIMARY KEY,
  order_id text NOT NULL, action text NOT NULL,
  authorization_id text NOT NULL, effect_input jsonb NOT NULL,
  state text NOT NULL CHECK(state IN ('awaiting_dispatch','dispatched')),
  origin text NOT NULL CHECK(origin IN ('consumer','legacy-import')),
  origin_evidence text NOT NULL,
  UNIQUE(order_id,action,authorization_id),
  FOREIGN KEY(order_id,action,authorization_id)
    REFERENCES authorizations(order_id,action,authorization_id)
);
\end{BookCode}

\code{processed\_events} 的联合主键仲裁同一 I。\code{message\_input} 保存完整原报文，\code{result} 保存这次处理的稳定结果，\code{fulfillment\_id} 引用真实 F。\code{pending} 仅用于受控事务内部构造，入口在提交以前必须完成终态或整笔回滚。下列 CHECK 约束状态和空值关系。提交终态的规则由受控入口兑现，数据库写权限限制到该入口。

\begin{BookCode}
CREATE TABLE processed_events (
  handler_scope text NOT NULL, source_id text NOT NULL,
  message_kind text NOT NULL, message_id text NOT NULL,
  message_input jsonb NOT NULL,
  status text NOT NULL CHECK(status IN ('pending','done','rejected')),
  fulfillment_id bigint REFERENCES fulfillments(fulfillment_id),
  result jsonb,
  PRIMARY KEY(handler_scope,source_id,message_kind,message_id),
  CHECK((status='pending' AND result IS NULL AND fulfillment_id IS NULL)
    OR (status IN ('done','rejected') AND result IS NOT NULL)),
  CHECK(handler_scope<>'fulfillment.initial.v1'
    OR status='pending'
    OR (status='done' AND fulfillment_id IS NOT NULL)
    OR (status='rejected' AND fulfillment_id IS NULL))
);
\end{BookCode}

入口逐项核对跨字段语义：接收结果中的 O、授权、F 与效果输入对应，F 的快照与授权行一致。外键检查被引用行及指定列存在，JSON 中的金额、商品等字段则由入口按保存版本核对。购买库的 O 属于另一数据库，履约库保存其受信标识和原事实，本地外键连接的是同库授权与 F。

\subsection*{先走通三张核心表之间的成功路径}

取履约库为空、EV1 的来源与输入已经核验、初次履约允许建立的共同初态。Q 先在 I1 的唯一键处取得这份接收身份，再按 B1 锁定授权行。B1 尚无 F 时插入履约单，设实际返回编号为 501，记作 F1；然后把包含 501 的 \code{created} 结果与 I1 的终态共同保存。另一个事件 I2 指向同一 B1 时，读取并比较 F1 的完整效果输入，保存引用 501 的 \code{existing} 结果。

图~\ref{fig:ch10-core-q} 只展示这三个对象之间的原子关系。完整 Q 还要接住两种继续责任：成功以后向购买库回报，发生输入冲突时留下可处置事实。下面依次补齐这两层，再展开包含成功、拒绝与异常的完整代码。

\CTenFigure{ch10-core-q}{I 仲裁本次接收，B 仲裁初次履约，结果引用实际返回的 F}{fig:ch10-core-q}

\subsection*{第二层：把成功后的回报保存为待办}

新 F 的建立还要留下一项回报购买库的持久工作。回报种类为 \code{fulfillment.accepted}，表示履约库已经受理建单。每个 F 只产生一项这种回报。待发内容可从同事务保存的 F 及其不可变效果输入取得。调度与受理证据沿第九章 Outbox 的规则另外保存，\code{ready/accepted} 在这里仅列出当前需要的两种阶段。第~\ref{section:10-4}~节再把它接到购买库的履约进度和实际发货。

\begin{BookCode}
CREATE TABLE fulfillment_reports (
  report_id bigint GENERATED BY DEFAULT AS IDENTITY PRIMARY KEY,
  fulfillment_id bigint NOT NULL UNIQUE REFERENCES fulfillments(fulfillment_id),
  order_id text NOT NULL,
  report_type text NOT NULL CHECK(report_type='fulfillment.accepted'),
  phase text NOT NULL CHECK(phase IN ('ready','accepted'))
);
\end{BookCode}

\subsection*{第三层：把暂时不能接受的输入交给核查}

需要隔离的输入也要有持久位置。例如 EV1 的完整输入与原绑定矛盾，自动建单无法替核查者判断哪份事实正确。\code{consumption\_cases} 保存完整来件、原绑定、原因、负责角色和决定。开放项保留在可扫描范围，经授权处置以后，决定与重新安排或导入效果共同保存。重复同一冲突按完整结构化来件恢复同一核查项。

\begin{BookCode}
CREATE TABLE consumption_cases (
  case_id bigint GENERATED BY DEFAULT AS IDENTITY PRIMARY KEY,
  handler_scope text NOT NULL, source_id text NOT NULL,
  message_kind text NOT NULL, message_id text NOT NULL,
  reason text NOT NULL, incoming jsonb NOT NULL, original jsonb,
  owner_role text NOT NULL,
  state text NOT NULL CHECK(state IN ('open','resolved')),
  decision jsonb,
  UNIQUE(handler_scope,source_id,message_kind,message_id,reason,incoming),
  CHECK((state='open' AND decision IS NULL)
     OR (state='resolved' AND decision IS NOT NULL))
);
\end{BookCode}

这份定义保留完整 JSON 以便逐项判断。体积很大的消息若要用摘要加速定位，可以先按摘要找到候选，再比较另存的完整内容。若把摘要相同直接视为内容相同，就要另外承担摘要碰撞的假设。相同核查输入过大而不适合普通索引时，受控存储也要保留原件，并为同一完整来件提供唯一仲裁。

把这五张表放在一起，可以看到同库外键检查的范围。接收记录和回报引用真实 F，F 的 B 引用授权行；核查项保存来件 I 以便定位，自动建单暂停时也能独立保留完整责任。

\CTenFigure{ch10-consumer-schema}{同库外键连接授权、履约、接收结果与回报，核查表独立保存未完成责任}{fig:ch10-consumer-schema}

\subsection*{把三层关系放回完整接收事务}

现在从履约库空初态执行 EV1，采用 PostgreSQL 18 的 Read Committed 和已确认提交持续保留的前提~\cite{pg}。接收事务先尝试插入 I1，用 \code{ON CONFLICT DO NOTHING RETURNING message\_id} 取得本次是否占有身份的真实结果。返回零行时，另用一条语句读已提交记录。Read Committed 为新语句取得新的查询快照，这时可以读到刚刚结束等待的竞争事务所提交的终态。

首次占有者随后按 B1 插入或查找授权行并加锁。来自既定成功事件的新授权按既定业务规则判断，迁移预先登记的 \code{hold} 和拒绝决定保留。读到相同 B 却有不同效果输入时，停止自动建单，保存原输入与来件的冲突。

允许建单时，插入 F 并取得 \code{RETURNING fulfillment\_id} 的真实值。以下以返回 F1 叙述，后续引用都沿这个实际返回值传递。另一事件若已建立同 B 的 F，插入返回零行，随后读取原 F，逐字段核对原效果输入，再保存本次的 \code{existing} 结果。读不到原 F、归属不符或快照矛盾都有异常出口。

\begin{BookCode}
验证来源、结构及类型，按保存格式取得 I、完整输入和效果输入
开启履约库事务 Q
插入 I.pending，保存完整输入，ON CONFLICT DO NOTHING RETURNING 消息号
若返回零行：
    用新语句读原 I，缺失则回滚后重读或重试
    同完整输入且已终态：恢复保存结果，结束本地事务
    异输入：保留原 I/F，插入或恢复持久冲突核查，提交核查
若本次取得 I：
    按 B 插入或锁住实际授权行，核对效果输入与当前决定
    效果输入矛盾或 hold：
        删除本事务尚未提交的 I.pending，要求 RETURNING 一行
        共同保存完整来件及有归属的核查，处理责任转入核查
    已有稳定业务拒绝或首次决定拒绝：
        保存 outcome=rejected、O、动作、原授权与原因
        条件回填 I.rejected，F 引用为空，要求 RETURNING 一行
    允许：
        按 B 插入 F，ON CONFLICT(B) DO NOTHING RETURNING 实际 F
        返回一行：核对真实 F/O/授权/输入，建立唯一回报待办
        返回零行：读取同 B 的原 F，核对全部效果输入
        保存完整 created 或 existing 结果及真实 F 引用
        条件回填本事务的 I.pending 为 done，要求 RETURNING 一行
    所有预期写入及终态核对通过后，共同提交 Q
本地依据已确定以后，在原 channel 确认本次投递，multiple=false
\end{BookCode}

效果输入矛盾或授权处于 \code{hold} 时，删除的只是本事务尚未提交的 I 候选。这个 Q 提交的是完整来件与核查责任，尚未把它声明为已接受的业务输入。同一来件重投会再次经过 I/B 仲裁，再按核查表的完整输入唯一关系恢复原核查项。原来已经提交的 I 及其输入始终保留，异参来件只在核查表中另存。

另一种设计可以在 I 上增加持久 \code{check} 状态和核查关联，使重复来件直接恢复隔离进度。采用它时，要一并规定核查决定怎样重新安排业务、怎样保留原输入与核查履历，以及 I 从隔离到完成时返回什么结果。本章采用独立核查责任，下面沿这套已经定义的状态继续推演。对频繁重投的隔离输入，两种设计还应在实际负载下比较写入、锁竞争和版本回收成本。同事务插入后删除会留下需要回收的版本，UPDATE 同样有版本维护成本，具体收益由负载与 VACUUM 行为决定~\cite{pg}。

为展示实际保存的结果，设本次插入返回编号 501，正文记作 F1：

\begin{BookCode}
{
  "outcome": "created",
  "order_id": "O1",
  "action": "initial",
  "authorization_id": "initial:O1",
  "fulfillment_id": 501,
  "state": "awaiting_dispatch"
}
\end{BookCode}

另一个事件恢复同一 F 时，本次保存的 \code{outcome} 为 \code{existing}，其余建单事实相同。这里的 \code{state} 属于建单受理结果快照。F 后来成为 \code{dispatched}，重复 I1 恢复原 \code{created} 结果，查询当前履约进度则从 F 的现态读取。两类读取的目的不同。

\CTenFigure{ch10-result-current}{I1 恢复建单时的完整结果，当前履约查询沿真实编号读取 F1 的新进度}{fig:ch10-result-current}

\noindent\begin{minipage}{\linewidth}
再看一次独立的拒绝历史。该订单的原商品不在当前支持范围，首次业务决定拒绝建单，稳定结果没有 F 引用：

\begin{BookCode}
{
  "outcome": "rejected",
  "order_id": "O1",
  "action": "initial",
  "authorization_id": "initial:O1",
  "reason": "unsupported-product"
}
\end{BookCode}
\end{minipage}

原事件再次到达，先恢复这个结果。后来部署支持该商品时，入口也先读取原拒绝。确需重新授权，负责者保存新的业务决定与明确的执行身份，沿相应操作继续。

F 插入返回零行后查不到合法原 F，回报待办或终态回填等预期写入返回零行，或者主键、效果唯一、外键等约束发生无法归入正常恢复分支的异常，事务都整笔回滚。身份候选、授权候选、F 和回报共同撤销。异常事务结束以后按原 I/B 重读：已有完整提交就恢复，没有提交且前提满足才重试，无法解释的关联进入有归属的核查。数据库断连使提交结果未知时，恢复者连接原库，用原 I 查询完整终态及关联，再决定恢复还是沿同一 I/B 重试。

\CTenFigure{ch10-commit-ack}{Q 的本地原子提交先确定恢复依据，随后 ACK 移交当前投递责任}{fig:ch10-commit-ack}

\subsection*{两个副本同时到达时}

W1、W2 同时处理相同 EV1。W1 插入 I1 后暂停，W2 也尝试插入 I1。唯一约束把二者放在同一仲裁位置，W2 等待 W1 的事务结局。W1 共同提交 I1/F1/回报，W2 随后走零行分支，用新语句看到原终态并恢复结果。两者都可以确认各自投递，新增 F 只有一份。

若 W1 回滚，W2 可以取得 I1 并成为首次提交者。再考虑调用方错误地混入不同内容：W1 带原 3000 分内容的候选回滚，W2 带同 EV1 的另一份 3500 分内容先提交，W1 恢复时便发现内容冲突。首次提交绑定的方案以先提交者确定原输入，先前回滚的候选已经撤销。可信发布者持续发送自己保存的原事件。消费者则用唯一约束仲裁提交，在已提交绑定上比较完整内容，并把异参输入交给持久核查。

\CTenFigure{ch10-original03}{相同接收身份的并发副本在唯一约束处等待并恢复原结果}{fig:ch10-concurrent}

上述迁移可以逐项检查不变量。每个 I 至多一份已提交绑定，绑定的完整输入及终态结果保持。每个 B 至多一份 F。履约范围内的成功接收结果引用同 B、同效果输入的真实 F，拒绝接收保存原因，独立审计范围引用自己的审计效果。已导入旧 F 具有迁移保存的来源依据。每个新建或导入的 F 同时具有一项回报依据。

初态为空，关系成立。首次成功在同一 Q 中同时建立 F、回报和终态。回滚保留此前已提交状态，已有身份恢复只读取原绑定。不同事件恢复原 F 时先比较效果输入，拒绝保存独立结论，隔离把未完成责任交给核查并保留原绑定。唯一约束、事务原子性和受控写入口共同支撑这组迁移，后续修改 F 的调度状态也保留这些身份与原结果。
